\documentclass[runningheads]{llncs}
\usepackage[english]{babel}
\usepackage[utf8]{inputenc}
% Carácteres especiales
\usepackage{fontenc}
% More advance packages

\usepackage{amssymb,amsbsy,mathrsfs,rotating,float,multirow,upgreek,amsmath,lineno,bm}

\usepackage{listings}
\usepackage{pgfplots}
\usepackage{tikz}
\usepgfplotslibrary{groupplots}
\usepackage{longtable}
\setlength{\LTcapwidth}{6in}
\usepackage[utf8]{inputenc}
\usepackage{epsfig,epic,eepic,threeparttable,amscd,here,lscape,tabularx,graphicx
}
\usepackage{booktabs}
\usepackage{tabu,array}
\usepackage{overpic}
\usepackage{microtype}
%%%%%%%%%%%
%\usepackage{newtxtext,newtxmath}

%%%%%%%%%%%
\usepackage{caption}
\usepackage{subcaption}
% Permite ver y configurar los parámetros de la página
\usepackage{layout}
%Hyperref permite ver las secciones del texto
\usepackage[hidelinks]{hyperref}
\usepackage[noabbrev,capitalize]{cleveref}

\usepackage{nomencl}

\providecommand{\ppunto}[2]{\langle#1, #2 \rangle}% producto punto
\providecommand{\promed}[1]{{\mathbb{E}}\left\lbrace
#1\right\rbrace}% operador de promedio
\providecommand{\promeddd}[2]{\mathbb{E}_{#1}\!\left\{#2\right\}}%
% op{\tiny }erador de promedio
\providecommand{\promedd}[2]{\mathds{E}_{#1}\left\{#2\right\}}%
% operador de promedio
\providecommand{\cov}[2]{{\fam=7 cov}\{#1, #2\}}% operador de covarianza
\providecommand{\gaus}[2]{{\fam=6 N}(#1 #2)} % FDP Gauss
% Kronecker delta function
\providecommand{\Kronecker}[2]{\delta_{K}( #1 , #2 )}
% cardinality
\providecommand{\Cardinality}[1]{| #1 |}

%%%% custom colors
\DeclareRobustCommand{\legend}[1]{%
  \textcolor{#1}{\rule{1ex}{2ex}}%
}
\colorlet{color0}{blue!40!gray}% pearson *
\colorlet{color1}{blue!80!gray} %pearson **
\colorlet{color2}{red!40!gray} %n-pearson*
\colorlet{color3}{red!80!gray}%n-pearson**
\colorlet{color4}{green!40!gray} %plv*
\colorlet{color5}{green!80!gray}

\definecolor{GII}{RGB}{255,255,107}
\definecolor{GI}{RGB}{72,197,72}
\definecolor{GIII}{RGB}{214, 36, 47}

\definecolor{frontal}{RGB}{141, 211, 199}
\definecolor{frontal_left}{RGB}{255, 237, 111}
\definecolor{frontal_right}{RGB}{255, 255, 179}
\definecolor{central_right}{RGB}{251, 128, 114}
\definecolor{central_left}{RGB}{204, 235, 197}
\definecolor{posterior_right}{RGB}{253, 180, 98}
\definecolor{posterior_left}{RGB}{217, 217, 217}
\definecolor{posterior}{RGB}{179, 222, 105}
\definecolor{crimson2143940}{RGB}{214,39,40}

%%%%
\definecolor{darkorange25512714}{RGB}{255,127,14}
\definecolor{forestgreen4416044}{RGB}{44,160,44}
\definecolor{steelblue31119180}{RGB}{31,119,180}

\providecommand{\est}[1]{{\widetilde {#1}}}
\providecommand{\s}[1]{\negthickspace#1\negthickspace}%
\newcolumntype{C}[1]{>{\centering}m{#1}}
\newcommand{\subconj}{\negthinspace\subset\negthinspace }
\newcommand{\en}{\negthinspace\in\negthinspace }
\newcommand{\igual}{\negthinspace=\negthinspace}
\newcommand{\dos}{\negthinspace:\negthinspace}

\newcommand{\Real}{\mathbb{R}}
\newcommand{\Natural}{\mathbb{N}}

\newcommand{\Tr}[1]{Tr \left( #1 \right)}
\newcommand{\diag}[1]{diag \left( #1 \right)}

\newcommand{\func}[1]{\mathcal {#1}}

\newcommand{\ve}[1]{\bm {#1}}
\newcommand{\mat}[1]{\bm {#1}}

\newenvironment{reviewer}{\setcounter{pointcounter}{1}}{}
\newcommand{\changes}[1]{\textcolor[rgb]{1.00,0.00,0.00}{#1}}
\newcommand{\genaro}[1]{\textcolor[rgb]{0.00,0.00,1.00}{#1}}
\newcommand{\point}[1]{\medskip \noindent
  \textsl{{\fontseries{b}\selectfont Point \thepointcounter}.
\stepcounter{pointcounter} #1}}
\newcommand{\reply}{\medskip \noindent \textbf{Response}.\ }
\newcommand{\initresponses}{\newcounter{pointcounter}}
\newcommand*{\captionsource}[2]{%
  \caption[{#1}]{%
    #1%
    \\\hspace{\linewidth}%
    \textbf{Source:} #2%
  }%
}

\providecommand{\am}[1]{\textcolor{green}{#1}}
\providecommand{\dg}[1]{\textcolor{blue}{#1}}
%
\begin{document}

We would like to extend our heartfelt thanks to each of you for your
time and effort in reviewing the document. Your input has been
invaluable in refining and enhancing the quality of the work. Below,
you will find our point-by-point responses to each of your comments.\\

\initresponses

\begin{reviewer}
  \textbf{Dear Dr. Cristian Alfonso Jiménez Castaño}\\

  \textbf{We sincerely thank you for all the valuable comments and
    recommendations you have provided. Your observations have been
  fundamental for improving the quality and clarity of the work.}

  \point{Document cohesion from the problem statement onwards.}

  \reply{The document has been restructured to improve
    cohesion between sections, especially in the methodology and
  presentation of results.}

  \point{Lack of bibliographic support in several
  subsubsections of the state of the art.}

  \reply{Appropriate bibliographic references have been
  added in all subsections that lacked support.}

  \point{Absence of justification for the use of
  convolutional networks vs. other architectures.}

  \reply{The methodological section has been expanded to
    include detailed justification for the choice of U-Net architecture
  and comparison with alternatives.}

  \point{Results and conclusions presented in
  separate and disconnected chapters.}

  \reply{The results sections have been unified and
  the connection with the research objectives has been improved.}

  \point{Figure 1-1: font size not recognizable at 100\%.}

  \reply{The font size has been increased and
  the figure has been resized for better visibility.}

  \point{Figure 1-2: small letters and page organization.}

  \reply{The figure has been resized and the
  layout has been reorganized to avoid blank spaces.}

  \point{Lines 429 and 440: usability frequency percentage not clear.}

  \reply{The methodology and study universe have been
    clarified for the calculation of the counted referenced  terms from
  the sources.}

  \point{Figure 1-3: font size and accommodation.}

  \reply{The font size has been increased and the
  distribution on the page has been improved.}

  \point{Excessive concentration on multiple annotators
  without considering other problems.}

  \reply{The section has been expanded to include a
  broader view of the state-of-the-art problems.}

  \point{Research question without clear methodology
  to justify U-Net.}

  \reply{The methodological justification
  for the choice of U-Net has been developed.}

  \point{Quick dismissal of other image types without
  justification.}

  \reply{Detailed justification has been added for the
  focus on histological images.}

  \point{General objective does not begin with verb in infinitive.}

  \reply{The general objective has been reformulated to
  begin with a verb in infinitive.}

  \point{Line 791: interpretability not reflected in
  specific objectives.}

  \reply{The objecives are no longer explicitly focused in the
    interpretability of the model, even though there is indeed a degree
  of interpretability.}

  \point{Line 792: very strong comparison that requires
  extensive experimental framework.}

  \reply{The scope of the comparison has been adjusted to
    be realistic with the proposed experimental framework, focused on
  the development of a reliability map with a satisfying performance.}

  \point{Section advances too much without presenting methodologies.}

  \reply{This section has been moved to the end of the document,
  after the conclusions so that the narrative flow is improved.}

  \point{Figure 2-2: improve resolution and size.}

  \reply{The resolution has been improved and the figure
  has been resized.}

  \point{Figure 2-3: remove black part of the interface.}

  \reply{The image has been cleaned to show only the
  pathological image.}

  \point{Section 2.2.2: too empty.}

  \reply{It has been completed with visual examples and
  detailed explanations.}

  \point{Section 2.3.1: inconsistency between labels and icons.}

  \reply{Inconsistencies have been corrected and
  the connection between paradigms has been improved.}

  \point{Line 970: clarify that they are machine learning
  paradigms in general.}

  \reply{The scope of the mentioned paradigms has been clarified.}

  \point{Subsection 2.3: too brief.}

  \reply{It has been expanded with additional development and
  bibliographic references.}

  \point{Subsections 3.1.2 and 3.1.3: without bibliographic support.}

  \reply{Appropriate bibliographic references have been added.}

  \point{Incorrect order of subsections.}

  \reply{It has been reordered to follow the correct logical sequence.}

  \point{Subsection 3.3: contribution not clear.}

  \reply{The purpose and contribution of
  this section have been clarified.}

  \point{Table 3-1: organize for better visibility.}

  \reply{The table has been reorganized to improve readability.}

  \point{Equation 4-10 too long.}

  \reply{It has been divided into intermediate functions for
  better readability.}

  \point{Section 4.1.4: this section should not be there without experiments.}

  \reply{It has been merged with the previous section for
  better coherence.}

  \point{Numbering begins at zero.}

  \reply{The numbering of the subsections of Chapter 5 have been
  corrected to begin at one.}

  \point{Figures 5-1 and 5-2 seem identical.}

  \reply{The duplicate figures have been unified.}

  \point{Lines 1627-1629: it is not clear if the
  regularization terms are scalars.}

  \reply{The nature of the regularization terms
  in equation 5-1 has been clarified.}

\end{reviewer}

\begin{reviewer}
  \textbf{Dear Dr. Juan Bernardo Gómez Mendoza}\\

  \textbf{We sincerely thank you for your positive evaluation and all
    the valuable
    comments and recommendations you have provided. Your observations
  have been crucial for improving the quality and completeness of the work.}

  \point{Lack of specific tests to illustrate the
  effect of each proposed modification.}

  \reply{Ablative experiments have been added to
  demonstrate the individual impact of each component of the solution.}

  \point{Insufficient characterization of the
  mathematical implications of the cost function.}

  \reply{The mathematical section has been expanded to
  include detailed analysis of the cost function properties.}

  \point{Lack of clarity in hyperparameter selection.}

  \reply{Optuna was used to select the hyperparameters through all
    experiments. The scalar experiment in particular was useful for
    determining the most meaningful hyperparameters for the model with
  the help of the Optuna B.}

  \point{Improvements in the visual presentation of
  figures and tables.}

  \reply{All figures and tables have been reviewed and
  improved for greater clarity.}

  \point{Consistency in mathematical notation.}

  \reply{All mathematical notation in the document has been
  reviewed and standardized.}

  \point{Improvements in writing and text clarity.}

  \reply{The writing throughout the document has been
  reviewed and improved for greater clarity.}

\end{reviewer}

\section*{Conclusion}

We have systematically addressed all the comments and
recommendations provided by both reviewers. The document has
been significantly improved in terms of:

\begin{itemize}
  \item Structural cohesion and logical flow
  \item Bibliographic support and references
  \item Clarity in methodology and justification of decisions
  \item Presentation of results and analysis
  \item Visual quality of figures and tables
  \item Consistency in notation and writing
\end{itemize}

We thank you again for your valuable time and constructive feedback,
which have been fundamental for elevating the quality of
this research work.

\end{document}
